\section{Approximation in a prescribed power coset}
\label{sec:approximation}

The closure theorem of Section~\ref{sec:arithmetic} permits approximation
at finitely many places within a prescribed coset of the subgroup generated
by powers.  We shall also need to control a local property at every other
place.  The construction in this section achieves that control when its
failure at integral points has geometric codimension at least two, provided
certain explicitly described nonintegral points satisfy the property as
well.  The output remains in the prescribed abstract coset throughout.

We adapt the line-avoidance and power-factor constructions of
\cite[Sections~2.3--2.4]{MPNF}.  In positive characteristic, a power map
need not have a nonzero differential.  We separate its inseparable part
before making the geometric codimension argument.  We also keep track of
the order in which the finite sets of exceptional places are chosen;
this will allow the construction to be applied when the local groups
arising later depend on the rational point being constructed.

\subsection{Avoidance in affine space}

For a place $v$, write $\mathcal O_v$ for the valuation ring of $k_v$,
$\kappa_v$ for its residue field, and $q_v=|\kappa_v|$.
An integral model of an affine $k$-variety means a model over the ring
of integers outside a finite set of places.  All assertions about
reduction below refer to fixed such models.  Enlarging the finite set
allows us to extend any fixed finite collection of morphisms and
equations to these models.

The following elementary form of avoidance leaves one place unrestricted.
That freedom will later be removed using a symmetry of the parameter map.

\begin{lemma}[Avoidance along a line]\label{approx:line-sieve}
Let $E$ be a finite-dimensional $k$-vector space and let $Z\subset E$
be a closed subset of geometric codimension at least two.  Choose
$d\in E(k)\setminus\{0\}$ such that its point $[d]$ at infinity
does not belong to the closure of $Z$ in $\bbP(E\oplus k)$.
Let
\[
 \pi:E\longrightarrow E/kd
\]
be the quotient map, and fix a rational linear section $s$ of $\pi$.
Fix a place $v_0$.

There are a finite set of places $\Sigma_0$ containing $v_0$ and an
integer $D\geq1$ with the following property.  Let $\Sigma\supseteq
\Sigma_0$ be finite, and prescribe nonempty open subsets
$\mathcal W_v\subset E(k_v)$ for $v\in\Sigma\setminus\{v_0\}$.
There exist $y\in E(k)$ and elements $l\in k$ of arbitrarily
large $v_0$-absolute value such that
\begin{enumerate}
\item $y+ld\in\mathcal W_v$ for every
      $v\in\Sigma\setminus\{v_0\}$;
\item $y+ld$ is integral outside $\Sigma$ and its reduction avoids
      $Z$ at every place outside $\Sigma$.
\end{enumerate}
One may take $\Sigma_0$ to consist of model exceptions, $v_0$, and
the places with $q_v\leq D$, where $D$ bounds the cardinalities of
the geometric fibers of the reductions of $\pi|_Z$.
\end{lemma}

\begin{proof}
The restriction $\pi|_Z$ is finite.  Indeed, projection from $[d]$
extends to a morphism on the projective closure of $Z$, since its
center is disjoint from that closure.  A fiber over an affine point
can acquire no point at infinity: its projective fiber line meets
the hyperplane at infinity only at $[d]$.  Thus $\pi|_Z$ is proper.
It has finite fibers, because a positive-dimensional closed subset
of a fiber line would be the whole line and would have $[d]$ in its
projective closure.  A proper morphism with finite fibers is finite.
Its image $Y\subset E/kd$ is closed, and
\[
 \dim Y\leq \dim Z\leq \dim E-2<\dim(E/kd),
\]
so $Y$ is proper.

After excluding finitely many places, the decomposition
$E=kd\oplus s(E/kd)$ is integral, $d$ is a basis of its first
summand, and $\pi|_Z$ extends to a finite morphism of the models.
A finite set of module generators for its coordinate algebra gives
a uniform bound $D$ for the cardinality of every geometric fiber.
Include these exceptional places and all places with $q_v\leq D$
in $\Sigma_0$.

Linear projection is open over each completion.  Additive strong
approximation away from $v_0$, equivalently the usual consequence
of Riemann--Roch for a function field, therefore gives
\[
 q\in (E/kd)(k)\setminus Y(k)
\]
that is integral outside $\Sigma$ and belongs to
$\pi(\mathcal W_v)$ at every prescribed place.  To ensure
$q\notin Y$, shrink one projected opening to avoid $Y$; a proper
algebraic subset has empty interior over an infinite local field.
If there are no prescribed openings, impose such an additional
opening inside the integral points at one auxiliary place.  This
does not change the integrality requirement.  Put $y=s(q)$.

The reduction of $q$ belongs to the reduction of $Y$ at only
finitely many places outside $\Sigma$.  For example, choose a
polynomial vanishing on $Y$ but not at $q$, clear its finitely many
denominators, and use the fact that its nonzero value at $q$ has
only finitely many zeros.  At each of these places the forbidden
values of the residue of $l$ number at most $D<q_v$.  There is
therefore a nonempty congruence condition on $l\in\mathcal O_v$
for which the reduction of $y+ld$ avoids $Z$.

At a prescribed place, $q\in\pi(\mathcal W_v)$ says precisely
that the set of $l\in k_v$ with $y+ld\in\mathcal W_v$ is
nonempty and open.  Additive strong approximation away from $v_0$
now supplies $l\in k$ satisfying these finitely many conditions
and integral at every remaining place outside $\Sigma$.
Shrink the openings for $l$ to be relatively compact.
There are infinitely many such $l$, since a nonempty open subset
of the adeles away from $v_0$ meets the dense diagonal copy of $k$
in infinitely many points.  If their absolute values at $v_0$
were bounded, these points would lie in a compact subset of the
full adele ring.  Its intersection with the discrete subgroup $k$
is finite.  Hence $|l|_{v_0}$ is unbounded, as required.
\end{proof}

\subsection{Products of torus powers}

For the remainder of this section, let $S=\SU(D,*)$ be
$k$-anisotropic of degree $n\geq3$, with notation as in
Section~\ref{sec:arithmetic}.  Fix $e\geq1$ and let
$R=S(k)^e$ denote the subgroup generated by all $e$-th powers.
We allow the ambient variety to be $S$ itself or a rational
translate $X=x_*S$ in $U$, where $x_*\in U(k)$.

Choose a maximal $k$-torus $T\subset S$ and a finite Galois
extension $M/k$ splitting $T$ and any fixed auxiliary data needed
in an application.  Write $p=\operatorname{char}k$.
Over $M$, identify $T$ with the determinant-one diagonal torus
in $\SL_n$.  Choose a cocharacter
\[
 \lambda:\bbG_{m,M}\longrightarrow T_M,
 \qquad
 \lambda(t)=\diag(t^{w_1},\ldots,t^{w_n}),
\]
where the integers $w_i$ are distinct, have sum zero, and are
not all congruent modulo $p$.  Such a choice exists for $n\geq3$:
start with a vector in $\bbZ^n$ of sum zero whose residues are
not all equal, and add multiples of $p$ of sum zero to make
its entries distinct.  In particular $d\lambda(1)$ is noncentral
in $\mathfrak{sl}_n$, even if $p$ divides $n$.

The restriction of scalars of $\lambda$, followed by the torus
norm, defines a $k$-morphism
\begin{equation}\label{approx:eta}
 \eta:\Res_{M/k}\bbG_m\longrightarrow T,
 \qquad
 \eta=N_{M/k}\circ\Res_{M/k}(\lambda).
\end{equation}
Over $M$ the coordinates of its source are indexed by
$\sigma\in\Gal(M/k)$, and its formula is
$\eta((a_\sigma)_\sigma)=\prod_\sigma
\lambda^\sigma(a_\sigma)$.
On the scalar copy of $\bbG_m$, this map has cocharacter
$\sum_\sigma\lambda^\sigma$.  That cocharacter is defined over
$k$ and must vanish, since $S$ is anisotropic.  Consequently
\begin{equation}\label{approx:scalar-invariance}
 \eta(ca)=\eta(a)
 \qquad(c\in K^\times,\quad a\in(M\otimes_k K)^\times)
\end{equation}
for every field extension $K/k$.

For $x\in X(k)$ and $k_1,\ldots,k_r\in S(k)$, consider
\begin{equation}\label{approx:product}
 \Phi(a_1,\ldots,a_r)
   =x\prod_{j=1}^r k_j\eta(a_j)^e k_j^{-1}.
\end{equation}
The parameter domain is the open torus
$\mathcal T_r=(\Res_{M/k}\bbG_m)^r$ in the affine space
$E_r=(\Res_{M/k}\bbA^1)^r$.  For rational parameters,
\begin{equation}\label{approx:exact-coset}
                    \Phi(\mathcal T_r(k))\subset xR.
\end{equation}
Indeed each factor is the $e$-th power of
$k_j\eta(a_j)k_j^{-1}\in S(k)$.

\begin{lemma}[Two independent parameter lists]\label{approx:two-lists}
Write $e=p^a e_0$ with $(e_0,p)=1$.  For a sufficiently long
list of rational conjugators $k_j$, the product map formed with
$\eta^{e_0}$ is smooth at the identity parameter tuple.
The conjugator tuples having this property form a nonempty
Zariski open subset of the corresponding power of $S$.

Choose two disjoint consecutive lists of factors with this
property.  In the full product map \eqref{approx:product}, allow
any further factors before or after these lists.  If
$Y\subset X$ is closed of geometric codimension at least two,
then $\Phi^{-1}(Y)\subset\mathcal T_r$ and its closure in $E_r$
have geometric codimension at least two.
\end{lemma}

\begin{proof}
Work first over an algebraic closure.  Put
$H=d\lambda(1)$, a noncentral diagonal trace-zero matrix.
The span of its conjugates contains every root vector $E_{ij}$.
Indeed conjugation by $1+tE_{ij}$ subtracts from $H$ a nonzero
multiple of $tE_{ij}$ whenever its $i$-th and $j$-th diagonal
entries differ, and permutation conjugations supply every pair
of positions.  The span is conjugation-invariant.  Conjugating
$E_{ij}$ by $1+tE_{ji}$ and subtracting root-vector terms
therefore supplies $E_{ii}-E_{jj}$.  These root vectors and
diagonal differences span $\mathfrak{sl}_n$ in every
characteristic.  In particular this argument does not require
$\mathfrak{sl}_n$ to be a simple Lie algebra.

The differential of $\eta$ at $1$ contains the line spanned by
$H$: vary just the absolute coordinate indexed by the identity
of $\Gal(M/k)$.  Finitely many conjugates of that line span
$\Lie(S)$, and multiplication by $e_0$ does not change the span.
Thus suitable conjugators give a surjective differential for
the product map at the identity.  Surjectivity is an open
condition on the conjugators.  Since $S(k)$ is Zariski dense,
this nonempty open contains a rational tuple.  Both source and
target are smooth, so the differential criterion gives the
asserted smoothness.

Let $\Phi_0$ denote the product with exponent $e_0$ in place
of $e$.  Each chosen list has a proper closed nonsmooth locus
in its own parameter torus.  The full map $\Phi_0$ is smooth
wherever either list is smooth, since multiplication by the
remaining factors is a translation when their parameters are
fixed.  The simultaneous nonsmooth locus has codimension at
least two, because the two lists use disjoint coordinates.
On its complement, smooth pullback preserves the lower bound
of two for the codimension of $Y$.  Hence
$\Phi_0^{-1}(Y)$ has that bound everywhere.

Finally,
\[
 \Phi=\Phi_0\circ[p^a],
\]
where $[p^a]$ raises every parameter to its $p^a$-th power.
After splitting the parameter torus, this is coordinatewise
Frobenius.  It is a finite universal homeomorphism, and
therefore preserves the dimensions and codimensions of inverse
images of closed subsets.  This proves the assertion for
$\Phi^{-1}(Y)$.  Passing from an open subset of affine space
to its closure does not change its dimension.
\end{proof}

We call either of the two lists in Lemma~\ref{approx:two-lists}
a \emph{basic list}.  Its actual map with exponent $e$ is
dominant, although it need not be smooth.  We need one further
consequence of dominance, uniform in the translating point.

\begin{lemma}[Integral parameters on one list]
\label{approx:residue-list}
Fix a basic list and a proper closed subset $Y\subset X$.
Outside finitely many places, depending on this list and these
fixed models but not on $x$, the following holds: for every
integral $x\in X(k_v)$, there are unit parameters on that list
for which the reduction of its product with $x$ avoids $Y$.
All other factors may be given parameter $1$.
\end{lemma}

\begin{proof}
The surjective differential for the exponent-$e_0$ map remains
surjective on reduction outside finitely many places.  Its
reduction is dominant, as is the exponent-$e$ map obtained by
composing with the surjective power map on the parameter torus.
We also exclude places where the reduction of $Y$ is not proper.

Choose affine coordinates on the underlying parameter space
and equations for $Y$.  Pulling these equations back through
the product map, and clearing the fixed torus denominators,
gives polynomials of bounded degree.  Their degrees are
independent of the translating point $x$.  For every residue
point $\bar x\in X(\kappa_v)$, dominance implies that at least
one such polynomial is nonzero.  Multiply it by the product
of the norm polynomials defining the torus boundary.  The
result is a nonzero polynomial, of degree bounded by a fixed
integer $C$, whose nonvanishing gives unit parameters and
avoidance of $Y$.

A nonzero polynomial of degree at most $C$ in $N$ variables
over a field with $q$ elements has at most $Cq^{N-1}$ zeros;
this follows by induction on $N$.  It therefore has a
nonvanishing value when $q>C$.  Exclude the finitely many
places with $q_v\leq C$, and lift the resulting parameters
to local units.
\end{proof}

\subsection{Boundary conditions and the approximation theorem}

The boundary of $\mathcal T_r$ in $E_r$ is a union of
hyperplanes after extending scalars to $M$:
\[
 H_{j,\sigma}=\{a_{j,\sigma}=0\},
 \qquad 1\leq j\leq r,\quad \sigma\in\Gal(M/k).
\]
We call a reduction a \emph{single-pole reduction} if it lies
on exactly one of these hyperplanes.  At a place unramified
in $M$, Frobenius must fix the unique vanishing coordinate.
Its action on the embeddings of the Galois extension $M/k$
is regular; hence this place splits completely in $M$.
Near this boundary hyperplane, away from all the others,
\eqref{approx:product} has the form
\begin{equation}\label{approx:single-pole}
                 \Phi=g_L\lambda^\sigma(t)^e g_R,
\end{equation}
where $t=a_{j,\sigma}$ and the two side factors are regular.
At a single-pole reduction outside model exceptions, $t$ has
positive valuation and the side factors are integral and
invertible.  Thus the necessary local analysis can be done
in split matrices, even though $S$ is anisotropic over $k$.

In the theorem below, a \emph{pole test} means finitely many
algebraic nonvanishing conditions on
$H_{j,\sigma}\setminus\bigcup_{(i,\tau)\ne(j,\sigma)}H_{i,\tau}$.
The conditions may involve $x$ and the conjugators as
coefficients.  They must be taken with all their Galois
conjugates.  The theorem states explicitly the nonemptiness
requirement that applications must verify.

\begin{proposition}[Approximation by power factors]
\label{approx:power-factor}
Let $k$ be a global function field, let $S=\SU(D,*)$ be
$k$-anisotropic of degree $n\geq3$, and put $R=S(k)^e$ for
$e\geq1$.  Let $X=x_*S$ with $x_*\in U(k)$, fix
$\gamma\in S(k)$, and set
\[
                   \mathcal C=x_*\gamma R\subset X(k).
\]
Let $Y\subset X$ be closed of geometric codimension at least
two.  Fix the torus, Galois splitting field, cocharacter and
map $\eta$ of \eqref{approx:eta}.  Let $B$ be a finite set
containing $A$, the exceptions for the fixed models, and a
place $v_0$ at which $S$ is isotropic.

At $B$, prescribe nonempty local open subsets of $X(k_v)$
whose product meets the closure of $\mathcal C$ given by
Proposition~\ref{arith:closure}.  Let
$(\mathcal U_v)_{v\notin B}$ be a family of subsets
$\mathcal U_v\subset X(k_v)$ satisfying the following conditions:
\begin{enumerate}
\item For every $v\notin B$, each integral point whose
      reduction avoids $Y$ belongs to $\mathcal U_v$.
\item For every $v\notin B$, the set $\mathcal U_v$ contains
      a nonempty local open subset.
\item There are two disjoint consecutive basic lists and,
      for each number of appended factors, a fixed finite
      family of algebraic pole tests, independent of $v$.
      After the initial point and conjugators have been
      fixed, these tests imply membership in $\mathcal U_v$
      at every single-pole reduction outside one finite
      set of model exceptions.  The tests admit choices in
      the following order.  The conjugator tuple of the
      two basic lists may be chosen in a nonempty Zariski
      open subset.  For
      each such tuple, the initial point $x$ may be chosen
      in a nonempty Zariski open subset of $X$.  For every
      such $x$ and every prescribed finite number of
      appended factors, their conjugators may be chosen
      in a nonempty Zariski open subset of the corresponding
      power of $S$.  For each tuple so chosen, the tests
      define a nonempty open subset of every absolute
      boundary hyperplane away from the other hyperplanes.
      Setting appended parameters to $1$ preserves all
      previously available pole tests.
\end{enumerate}
Then there is $x'\in\mathcal C$ satisfying the prescribed
open conditions at $B$ and belonging to $\mathcal U_v$ for
every $v\notin B$.

Finitely many additional nonempty Zariski open conditions
may be imposed at the three stages in \textup{(iii)}, as
long as they remain compatible with the stipulated choices.
\end{proposition}

\begin{proof}
We shall first choose the group-valued product and then apply
Lemma~\ref{approx:line-sieve} to its affine parameters.  There
are two reasons to enlarge the initial finite set of places:
the initial rational point may have denominators, and the
eventual projection may introduce further model exceptions.
The former will be handled by extra factors; the latter by
the fixed basic lists.

\emph{Choosing the basic lists and the initial point.}
Choose the conjugators of the two basic lists in the open
subset supplied by \textup{(iii)}.  Nonempty Zariski opens
can be intersected here, since the product of copies of $S$
is geometrically irreducible.  Put into $B$ the model
exceptions for these fixed lists and the finite set from
Lemma~\ref{approx:residue-list}.  At each newly added place,
choose an opening contained in $\mathcal U_v$ using
\textup{(ii)}.  These places lie outside $A$, so the
additional conditions are compatible with the closure
of $\mathcal C$.

Choose $x\in\mathcal C$ meeting these local conditions and
the nonempty Zariski open condition on $x$ in
\textup{(iii)}.  This follows from
Proposition~\ref{arith:closure}: use an omitted isotropic
place different from the finitely many places at which
approximation is required.  The Zariski condition can be
included by shrinking the opening at $v_0$, because a
nonempty analytic open subset of a smooth geometrically
irreducible variety is Zariski dense.

Let $B_1$ be the finite set of places outside the enlarged
$B$ at which $x$ is not integral.  At each $v\in B_1$ choose
a nonempty local open subset of $\mathcal U_v$.

\emph{Moving the denominator places.}
For $v\in B_1$, the group $S$ is isotropic over $k_v$.
The image
\[
 \eta\bigl((M\otimes_k k_v)^\times\bigr)^e
       \subset S(k_v)
\]
is noncentral.  Indeed its geometric image contains the
noncentral cocharacter image, and local points of its
source torus are Zariski dense.  Local Kneser--Tits and
Tits simplicity, as recalled in Section~\ref{sec:arithmetic},
imply that its normal closure is all of $S(k_v)$.
To be explicit, simplicity modulo the center first gives
surjectivity onto that quotient; the remaining quotient
is abelian and vanishes because $S(k_v)$ is perfect.
Every element of this normal closure is a finite product
of conjugates of the displayed values; inverses are
obtained by inverting the parameters.

Consequently finitely many factors of the form in
\eqref{approx:product} can move $x_v$ into the chosen
opening at every $v\in B_1$.  Use the largest number needed
at these finitely many places, inserting identity factors
where necessary.  Approximate the local conjugators by
rational conjugators, using strong approximation for $S$
away from $v_0$~\cite{PrasadSA}; the hypothesis is satisfied because
$S$ is simply connected and $S(k_{v_0})$ is noncompact.
Choose the rational conjugators integral outside
$B\cup B_1$.  The nonempty Zariski open condition on the
appended conjugators in \textup{(iii)} can be imposed at
the same time: first shrink their local openings into
that Zariski open and then apply strong approximation.

With these conjugators fixed, continuity gives nonempty
open conditions on the parameter tuple at each place of
$B\cup B_1$.  At $B$ take all parameters close to $1$.
At $B_1$ take the basic parameters close to $1$ and the
appended parameters close to the values effecting the
local moves.  Let $r$ be the resulting total number of
factors and retain the notation $\Phi:\mathcal T_r\to X$.

\emph{The subset to be avoided.}
In $E_r$ form the union $Z$ of the following closed subsets:
the closure of $\Phi^{-1}(Y)$; every intersection of two
distinct absolute boundary hyperplanes; and, on each
boundary hyperplane, the closure of the locus where its
pole tests fail away from the other hyperplanes.  Include
all Galois conjugates, so this union is defined over $k$.
Lemma~\ref{approx:two-lists} gives codimension at least two
for the first part.  The pairwise intersections have
codimension two, and the remaining parts do also have
codimension at least two, since each test succeeds on a
nonempty open subset of its hyperplane.

Outside finitely many model exceptions, integral parameters
whose reduction avoids $Z$ have one of two forms.  If
every absolute coordinate is a unit, $\Phi$ is integral
and its reduction avoids $Y$, so \textup{(i)} applies.
Otherwise exactly one coordinate has positive valuation,
and its pole tests hold, so \textup{(iii)} applies.
In either case $\Phi$ belongs to $\mathcal U_v$.

\emph{Choosing the line and treating the remaining exceptions.}
Choose $d=(d_1,\ldots,d_r)\in E_r(k)$ whose point at
infinity misses the projective closure of $Z$, and with
each $d_j$ sufficiently close to $1$ at $v_0$ that
$\Phi(d)$ satisfies the prescribed opening there.
These requirements are compatible: the directions excluded
at infinity form a proper algebraic subset, whereas the
condition at $v_0$ is a nonempty analytic open.  Weak
approximation on affine space supplies such a rational
direction.  Choose a rational linear section of projection
along $d$.

The equations defining $Z$, the projection, and its fiber
bound now give a fixed finite set of additional exceptional
places as in Lemma~\ref{approx:line-sieve}.  Let $\Sigma$
contain those places, the exceptions in \textup{(iii)},
all exceptions to the preceding reduction argument, and
$B\cup B_1$.  At a place of
$\Sigma\setminus(B\cup B_1)$, the initial point and every
conjugator are integral.  Lemma~\ref{approx:residue-list}
therefore gives unit parameters on one basic list whose
product with $x$ reduces outside $Y$; set all other
parameters to $1$.  A sufficiently small local opening
around this tuple maps into $\mathcal U_v$ by
\textup{(i)}.  This argument is valid even when the
projection or the pole equations have unsuitable
reduction at that place.  The small residue fields for
which the basic-list argument could fail were already
included in $B$ before $x$ was chosen.

Thus there are parameter openings at every place of
$\Sigma$, with the required local consequence, and no
further choices of rational conjugators or initial points
are needed.  Apply Lemma~\ref{approx:line-sieve}, omitting
the opening at $v_0$.  It gives $y\in E_r(k)$ and
parameters $a=y+ld$ with $|l|_{v_0}$ arbitrarily large,
meeting all prescribed openings away from $v_0$ and
reducing outside $Z$ at every place outside $\Sigma$.

\emph{Recovering the opening at $v_0$ and the exact coset.}
For each component of the parameter tuple,
\[
 \eta(y_j+ld_j)
    =\eta(d_j+l^{-1}y_j)
    \longrightarrow\eta(d_j)
       \qquad\text{as }|l|_{v_0}\longrightarrow\infty,
\]
by \eqref{approx:scalar-invariance}.  Hence $\Phi(a)$
eventually satisfies the opening at $v_0$.  For such $l$
all components $a_j$ are nonzero at $v_0$, so they belong
to $M^\times$ and $a\in\mathcal T_r(k)$.
At the other places of $\Sigma$ the chosen openings give
the required conclusion, and outside $\Sigma$ avoidance
of $Z$ gives it.  Finally \eqref{approx:exact-coset} gives
\[
                       x'=\Phi(a)\in xR=\mathcal C.
\]
The conditions selected at the places added to $B$ were
contained in their original allowed sets $\mathcal U_v$,
so the conclusion holds for the original finite set as
well.
\end{proof}

\begin{remark}\label{approx:open-locus}
The proposition also applies to local properties that are
only formulated on a nonempty Zariski open $X^0\subset X$.
Include $X(k_v)\setminus X^0(k_v)$ in each allowed set,
and refine the opening at $v_0$ to lie in $X^0(k_{v_0})$.
Then the global point belongs to $X^0(k)$ and has the
asserted properties there.  This will be used for regular
semisimple elements and for a later open set on which
certain auxiliary algebras are defined.

In an application with prescribed determinant but no
prescribed abstract coset, first use weak approximation
on the determinant slice to meet the finitely many local
conditions and then choose the coset containing that
point.  Conversely, when a coset is prescribed, the proof
never replaces it by its closure: every change of the
initial rational point is multiplication by explicit
$e$-th powers.
\end{remark}
